\documentclass[10pt,a4paper]{amsart}
\usepackage[margin=19mm]{geometry}
\usepackage{amsmath,amssymb,mathtools}
\usepackage[scr=rsfs,cal=euler]{mathalfa}
\usepackage{xfrac}
\usepackage[unicode,hidelinks,bookmarks=false]{hyperref}
\newcommand{\ie}{i.e.\ }
\newcommand{\cf}{cf.\ }
\newcommand{\resp}{respectively\ }
\newcommand{\Subsection}{\subsection}
\providecommand{\nobreakdash}{\nobreak\hbox{-}}
\setlength{\emergencystretch}{2em}
\multlinegap0pt
\usepackage{amssymb}
\usepackage{mathtools}
\newtheorem{theorem}{Theorem}
\newcommand{\HH}{\mathbb H}
\DeclareMathOperator{\divisor}{div}
\title{A Fricke transformation for cubic-residue eta products of level 13}
\author{Cetin Hakimoglu-Brown}
\date{Source: arXiv:2609.06804v1 (CC BY 4.0)}
\begin{document}
\maketitle

\noindent\emph{Source and licence.} A Fricke transformation for cubic-residue eta products of level 13. Version arXiv:2609.06804v1 (CC BY 4.0), distributed by arXiv under the Creative Commons Attribution 4.0 International licence, http://creativecommons.org/licenses/by/4.0/, as recorded in the arXiv licence metadata for this exact version and shown on the abstract page https://arxiv.org/abs/2609.06804v1. Full licence notice: \texttt{licenses/arxiv-2609.06804-CC-BY-4.0.txt}.\par\smallskip

\noindent\textit{Editorial scope.} Counted unit: the article's theorem thm:main (source0.tex lines 73--84). The article's complete original proof of it is delivered with it (source0.tex lines 505--536, one \texttt{proof} environment of 32 lines, which the article places away from the statement). No mathematics is written, repaired or re-derived: the statement and the proof are the article's own lines, in the article's own notation, and no retained line is substituted or re-flowed.  Retained uncounted as the source-local context the proof needs: lines 233--235; lines 278--281; lines 308--310; lines 332--334; lines 440--443; lines 499--501.  Nothing was omitted from any retained span, so the package carries no omission record.  No retained line contains a citation, so the wrapper carries no bibliography and no bibliography entry.  No retained line contains a figure environment, an image-inclusion command or drawing code, so no image is delivered and the package declares no asset dependency.  Editorial formatting only, in addition: an AMS \texttt{amsart} preamble that restates the article's own declarations and macros used by the retained lines, namely the environments \texttt{theorem} on separate counters, together with the standard packages those lines need.  The retained environments print in this document as Theorem 1 in order of appearance, whereas the article prints the same items with the numbering of its own full document.  The complete native source of the article as distributed in the arXiv e-print for this exact version is bundled unchanged as \texttt{sources/209573/source0.tex}.
\begin{equation}\label{eq:normalization}
 G_0=q^{-1/6}P_0,\qquad G_1=q^{-1/6}P_1,\qquad G_2=q^{5/6}P_2.
\end{equation}

\begin{equation}\label{eq:abc}
 \alpha=\beta+\gamma,\quad \gamma^2+\alpha\beta=13,
 \quad \beta^2+\alpha\gamma=13,\quad \alpha\beta\gamma=13.
\end{equation}

\begin{equation}\label{eq:fricke-image}
 G_j(W\tau)=\frac{q^{-1/6}}{4s_jT_j(\tau)}.
\end{equation}

\begin{equation}\label{eq:gauss}
 4s_0\alpha=4s_1\beta=4s_2\gamma=\sqrt{13}.
\end{equation}

\begin{equation}\label{eq:divisors}
 \divisor(t)=[c_1]+[c_5]-D,\qquad
 \divisor(G_1^6)=-\mathcal C+6D.
\end{equation}

\begin{equation}\label{eq:projective}
 t(W\tau)=\frac{\gamma-\beta t(\tau)}{\beta+\alpha t(\tau)}.
\end{equation}

\begin{theorem}\label{thm:main}
For every $\tau\in\HH$, the vector
\[
 H(\tau)=\begin{pmatrix}q^{-1/6}P_0(q)\\q^{-1/6}P_1(q)\end{pmatrix}
\]
satisfies
\begin{equation}\label{eq:main}
 H\left(-\frac1{13\tau}\right)=S_{13}H(\tau),\qquad
 S_{13}=\frac1{\sqrt{13}}\begin{pmatrix}\gamma&\beta\\\alpha&-\gamma\end{pmatrix}.
\end{equation}
Moreover $S_{13}^2=I$ and $H(\tau+1)=e^{-\pi i/3}H(\tau)$.
\end{theorem}

\begin{proof}[Proof of Theorem~\ref{thm:main}]
The numerator and denominator in~\eqref{eq:projective} have no common zero, since $\beta^2+\alpha\gamma=13$. At a pole of $t$ their ratio has a finite value. The polar divisor of $t\circ W$ is $W(D)$, where the involution identifies pushforward and pullback. It follows that
\begin{equation}\label{eq:denominator-divisor}
 \divisor(\beta+\alpha t)=W(D)-D.
\end{equation}
In particular this function does not vanish on $\HH$.

Now define the nonvanishing holomorphic function
\[
 R(\tau)=\frac{\sqrt{13}\,G_1(W\tau)}{(\beta+\alpha t(\tau))G_1(\tau)}.
\]
Its sixth power is an invariant modular function on $\Gamma_1(13)$, by the modularity of $G_1^6$ and $t$, and the normalization of the group by $W$. Equations~\eqref{eq:divisors} and~\eqref{eq:denominator-divisor} give
\begin{align*}
 \divisor(R^6)
 &=(-\mathcal C+6W(D))-(-\mathcal C+6D)-6(W(D)-D)\\
 &=0.
\end{align*}
Hence $R^6$ is constant on $X_1(13)$. Because $R$ is holomorphic on the connected domain $\HH$, it is constant as well. As $\tau\to i\infty$,
\[
 G_1(W\tau)\sim\frac{\beta}{\sqrt{13}}q^{-1/6},\qquad
 G_1(\tau)\sim q^{-1/6},\qquad t(\tau)\to0,
\]
by~\eqref{eq:fricke-image} and~\eqref{eq:gauss}. Therefore $R=1$, and
\[
 G_1(W\tau)=\frac{\beta+\alpha t}{\sqrt{13}}G_1
 =\frac{\alpha G_0-\gamma G_1}{\sqrt{13}}.
\]
Using~\eqref{eq:projective} gives
$G_2(W\tau)=(\gamma-\beta t)G_1/\sqrt{13}$, and adding these two identities yields
$G_0(W\tau)=(\gamma G_0+\beta G_1)/\sqrt{13}$.
This is~\eqref{eq:main}. Finally~\eqref{eq:abc} gives $S_{13}^2=I$, and~\eqref{eq:normalization} gives the translation law.
\end{proof}

\end{document}
